% Section 3.11, from lectures/L03/appendix/b_exercises.md.

\section{Exercises}
\label{c3:sec:exercises}

\begin{aside}
\textbf{How to check your work.} Every exercise below that asks you to write a VHDL module comes
with a self-checking testbench under \file{lectures/L03/exercises/}. Write your module in its
directory \file{exercises/<module>/}, with the entity name and the \textbf{port order} the exercise
states, and then run it with GHDL \textemdash{} see \secref{c2:sec:testbenches} for the three
commands. This is the first chapter whose exercises need that for a clocked design, so it is worth
rereading before you start on \code{register4}.
\end{aside}

% ------------------------------------------------------------------------------------------
\exerciseset{The D latch}

\exercise{The first circuit that remembers}{Circuit}
\label{c3:ex:latch}
Build the first circuit in this book that can \emph{remember} anything.

Everything you have built so far has been combinational: the outputs follow the inputs, and the
moment the inputs change the outputs change with them. A latch is the smallest possible departure
from that. It has a data input \code{D} and a control input \code{enable}, and it has two outputs,
\code{Q} and its complement \code{Qn}. As long as \code{enable} is high the latch is
\textbf{transparent}: \code{Q} simply follows \code{D}, and it looks like a wire. What is
interesting is what happens when \code{enable} goes low. The latch \textbf{latches}, holding
whatever \code{D} was at that instant, and from then on \code{D} can do as it likes without
\code{Q} noticing. That held value is one bit of memory, and it is built from nothing but the gates
you already know, wired so that the output feeds back to the input.

\code{Qn} is not decoration. The feedback loop is what makes the circuit remember, and the two
outputs are what feed each other: \code{Q} is computed from \code{Qn} and \code{Qn} from \code{Q}.
Build it and you see the loop; that is the whole point of drawing it by hand before you even meet
the VHDL.

Use the equations from \secref{c3:sec:latch}:

\[ Q = (D' \cdot \mathit{enable} + Q_n)' \qquad Q_n = (D \cdot \mathit{enable} + Q)' \]

\xpart{a} Realise the corresponding gate network: draw the network by hand, and recreate and
simulate it in CircuitVerse.

\xpart{b} Test the transparent and the latched state. Set \code{enable = 1}, change \code{D}, and
observe what happens to \code{Q} and \code{Qn}. Then set \code{enable = 0}, change \code{D} again,
and observe what happens. Explain the difference between the two cases.

\xpart{c} Run \code{D} and \code{enable} through the combinations \code{00}, \code{01}, \code{10},
\code{11}, twice. Confirm that you observe both cases: the latch latching after being transparent
with \code{D = 1}, and the latch latching after being transparent with \code{D = 0}. Explain how
the value of \code{Q} depends on the value of \code{D} at the moment \code{enable} changes from
\code{1} to \code{0}.

% ------------------------------------------------------------------------------------------
\exerciseset{The D flip-flop}

\exercise{From latch to flip-flop}{Circuit}
\label{c3:ex:dff}
Extend the D latch from Exercise~\ref{c3:ex:latch} into a D flip-flop, following
\secref{c3:sec:dff}.

The flip-flop should have the inputs \code{clock}, \code{reset_n}, \code{D} and \code{enable}, and
the outputs \code{Q} and \code{Qn}.

The reset signal is \textbf{asynchronous} (it takes effect without waiting for a clock edge) and
\textbf{active low} (setting \code{reset_n = 0} should force \code{Q} to \code{0} immediately).

\task{Hints}
\begin{itemize}
  \item The reset has to do two things in each of the two internal latches, not one:
  \begin{itemize}
    \item Feed an active-high \code{reset} (\code{reset_n} through a NOT gate) into the NOR gate
      producing \code{Q}. That forces \code{Q} low without waiting for a clock edge.
    \item \code{AND} \code{reset_n} into the AND gate producing that latch's product of \code{D}
      and \code{enable}. That frees \code{Qn} to go high.
    \item Gating only the feedback path for \code{Q} is not enough: with \code{enable = 1} and
      \code{D = 1} it leaves \code{Qn} low, and a low \code{Qn} is exactly what holds \code{Q}
      high, so the reset would do nothing.
  \end{itemize}
  \item Drive the two internal latches' enable inputs with \code{clock} and the inverse of
    \code{clock} respectively.
  \item Follow \secref{c3:sec:dff} for how the external \code{enable} input controls whether the
    flip-flop takes a new value.
\end{itemize}

\xpart{a} Realise the corresponding gate network: draw it by hand, recreate and simulate it in
CircuitVerse, and set the clock period to \code{1000 ms}, slow enough to observe the behaviour by
eye.

\xpart{b} Test the clocked behaviour: change \code{D} during the clock's high phase, and then
during its low phase, and observe \code{Q} and \code{Qn}. Confirm that the outputs change only at
the active clock edge, never between clock edges. Explain how the two internal latches work
together to give edge-triggered behaviour.

\xpart{c} Test the asynchronous reset: first set \code{Q = 1}, then activate the reset by setting
\code{reset_n = 0}, and observe whether \code{Q} clears immediately or at the next clock edge.
Explain the result by reference to the difference between an asynchronous signal and a synchronous
one.

% ------------------------------------------------------------------------------------------
\exerciseset{Registers}

\exercise{\code{register4}}{VHDL}
\label{c3:ex:register4}
Write a complete VHDL entity named \code{register4}: four flip-flops side by side, sharing one
clock, one reset and one enable.

This is the exercise where the flip-flop stops being a curiosity and becomes a building block. A
lone D flip-flop stores one bit, which is not much use in itself; what a design actually wants is
somewhere to keep a \emph{number}. Put four flip-flops in a row, connect the same clock to all of
them, and you have a four-bit register. Nothing about the individual flip-flop changes. Only the
width does.

What is worth noticing is how little the VHDL has to change to say that. You write the same
synchronous process template from \secref{c3:sec:template} that you already used for one bit, with
\code{std_logic_vector(3 downto 0)} instead of \code{std_logic}. There is no loop, no repetition, no
four copies of anything. That is not a shortcut the language is offering you; it is a fact about the
hardware, and part \textbf{c)} below asks you to say why.

The entity should have:

\begin{booktable}{@{}p{5em}p{3.4em}p{10.8em}L@{}}
  \thead{Port} & \thead{Dir.} & \thead{Type} & \thead{Description} \\
  \midrule
  \code{clock} & in & \code{std_logic} & System clock. \\
  \code{reset_n} & in & \code{std_logic} & Asynchronous, active low. Clears the register to
    \code{"0000"}. \\
  \code{enable} & in & \code{std_logic} & \code{'0'} keeps the current value; \code{'1'} captures
    \code{d} at the next rising edge. \\
  \code{d} & in & \code{std_logic_vector(3 downto 0)} & The value to capture. \\
  \code{q} & out & \code{std_logic_vector(3 downto 0)} & The stored value. \\
\end{booktable}

Write the process out yourself rather than copying the worked flip-flop example in
\secref{c3:sec:dffvhdl}. Typing it in is what makes the template stop being something you recognise
and start being something you can produce from memory, and you will write it in every remaining
chapter of this book.

\xpart{a} Write the entity and the synchronous process, adapting the template from a single bit to
a four-bit vector.

\xpart{b} Check it with the testbench, as described at the top of this section.

\xpart{c} Explain why the same template scales from one flip-flop to four with no structural
change, while a four-bit \emph{counter} cannot simply be four one-bit counters side by side. What
separates the two cases?

\modulefig[0.6]{lectures/L03/appendix/images/register4.png}

\begin{selfcheck}
\textbf{Self-check.} Name your entity \code{register4}, with the ports \code{clock},
\code{reset_n}, \code{enable}, \code{d} (\code{std_logic_vector(3 downto 0)}) and \code{q}
(\code{std_logic_vector(3 downto 0)}), declared in that order; its testbench is in
\file{exercises/register4/}.
\end{selfcheck}

% ------------------------------------------------------------------------------------------
\exerciseset{Edge detection}

\exercise{The toggle circuit by hand}{Circuit}
\label{c3:ex:edgehand}
Design a circuit that toggles an LED once for every sampled rising edge of a button signal. Follow
\secref{c3:sec:edge} and \secref{c3:sec:ledtoggle}.

The circuit should have the inputs \code{button}, \code{clock} and \code{reset}, and the output
\code{led}.

In this exercise \code{reset} is active high, unlike the reset in the worked example. It is still
\textbf{asynchronous}, like every reset in this book: activate it and \code{led} clears
immediately, without waiting for a clock edge. In the template from \secref{c3:sec:template} that
means the reset branch stays outside the \code{rising_edge(clock)} test, and only the polarity of
the reset test changes.

\task{Hints}
Use two flip-flops. One stores the previous sampled value of \code{button}, and that value is used
to detect a rising edge. The other stores the LED's current state, and is toggled only when a
rising edge is detected.

\xpart{a} Realise the corresponding gate network: draw it by hand, recreate and simulate it in
CircuitVerse, and set the clock period to \code{1000 ms}.

\xpart{b} Test the edge detector: hold \code{button} high for several consecutive clock cycles,
confirm that \code{led} toggles exactly once, and that it does not toggle once per clock cycle.
Explain which stored signal stops the circuit from repeatedly detecting the same high level as a
new rising edge.

\exercise{\code{led_toggle_single}}{VHDL}
\label{c3:ex:ledsingle}
Implement the circuit from Exercise~\ref{c3:ex:edgehand} in VHDL.

\xpart{a} Create an entity named \code{led_toggle_single} with:

\begin{booktable}{@{}p{5em}p{3.4em}p{5.6em}L@{}}
  \thead{Port} & \thead{Dir.} & \thead{Type} & \thead{Description} \\
  \midrule
  \code{clock} & in & \code{std_logic} & System clock. \\
  \code{reset} & in & \code{std_logic} & Asynchronous, and \textbf{active high}, unlike every other
    module in this book. Activate it and \code{led} clears without waiting for a clock edge, so the
    reset branch belongs outside the \code{rising_edge(clock)} test. \\
  \code{button} & in & \code{std_logic} & Toggles \code{led} once per press. \\
  \code{led} & out & \code{std_logic} & The LED's state. \\
\end{booktable}

Note the single-bit types: the worked example uses \code{std_logic_vector} ports for two buttons and
two LEDs, while this one drives a single pair.

Do not copy \file{led_toggle.vhd} wholesale. Do this instead:
\begin{itemize}
  \item Derive the two clocked processes yourself: one stores the previous value of \code{button},
    the other stores and toggles the current value of \code{led}.
  \item Derive the edge-detection equation for a rising edge and an active-high button, that is,
    \secref{c3:sec:edge}'s $\mathit{edge} = \mathit{current} \cdot \mathit{previous}'$.
  \item Do not use the falling-edge equation from the worked example, which assumes an active-low
    button.
\end{itemize}

\xpart{b} Verify the design with its testbench.

The testbench drives the same sequence you would go through by hand on a board:
\begin{itemize}
  \item It activates \code{reset} and checks that \code{led} is cleared.
  \item It pulls \code{button} high and checks that \code{led} toggles exactly once.
  \item It \textbf{holds \code{button} high} for several more clock cycles and checks, after each
    of them, that \code{led} does \emph{not} toggle again.
  \item It releases and presses again, and checks that the second press toggles \code{led} once
    more.
  \item It activates \code{reset} with \code{led} lit and \textbf{with no clock edge afterwards},
    and checks that \code{led} clears anyway. That is the check that fails if you put the reset
    inside the \code{rising_edge(clock)} branch.
\end{itemize}

The third check is the one that fails if your edge detector is wrong. An \code{led} toggling every
clock cycle while the button is held down means you drove the toggle from \code{button}'s
\textbf{level} instead of from its edge. Reread the equation above: \code{edge} is high only when
\code{button} is high \emph{now} and its stored previous value is low.

\modulefig[0.56]{lectures/L03/appendix/images/led_toggle_single.png}

\begin{selfcheck}
\textbf{Self-check.} Name your entity \code{led_toggle_single}, with the ports \code{clock},
\code{reset}, \code{button} (in) and \code{led} (out), declared in that order; its testbench is in
\file{exercises/led_toggle_single/}. Note that \code{reset} is \textbf{active high} here, unlike in
the worked example.
\end{selfcheck}
